\documentclass[10pt,a4paper]{amsart}
\usepackage{amsmath,amssymb,amsthm,graphicx}
% For commutative diagrams we use Xy-pic
\usepackage[all]{xy}

% We use 2cell for 2-commutative diagrams.
\xyoption{2cell}
% Theorem environments.
%
\theoremstyle{plain}
\newtheorem{theorem}[subsection]{Theorem}
\newtheorem{proposition}[subsection]{Proposition}
\newtheorem{lemma}[subsection]{Lemma}

\theoremstyle{definition}
\newtheorem{definition}[subsection]{Definition}
\newtheorem{example}[subsection]{Example}
\newtheorem{exercise}[subsection]{Exercise}
\newtheorem{situation}[subsection]{Situation}

\theoremstyle{remark}
\newtheorem{remark}[subsection]{Remark}
\newtheorem{remarks}[subsection]{Remarks}

\numberwithin{equation}{subsection}

% Macros
%
\def\lim{\mathop{\mathrm{lim}}\nolimits}
\def\colim{\mathop{\mathrm{colim}}\nolimits}
\def\Spec{\mathop{\mathrm{Spec}}}
\def\Hom{\mathop{\mathrm{Hom}}\nolimits}
\def\Ext{\mathop{\mathrm{Ext}}\nolimits}
\def\SheafHom{\mathop{\mathcal{H}\!\mathit{om}}\nolimits}
\def\SheafExt{\mathop{\mathcal{E}\!\mathit{xt}}\nolimits}
\def\Sch{\mathit{Sch}}
\def\Mor{\mathop{\mathrm{Mor}}\nolimits}
\def\Ob{\mathop{\mathrm{Ob}}\nolimits}
\def\Sh{\mathop{\mathit{Sh}}\nolimits}
\def\NL{\mathop{N\!L}\nolimits}
\def\CH{\mathop{\mathrm{CH}}\nolimits}
\def\proetale{{pro\text{-}\acute{e}tale}}
\def\etale{{\acute{e}tale}}
\def\QCoh{\mathit{QCoh}}
\def\Ker{\mathop{\mathrm{Ker}}}
\def\Im{\mathop{\mathrm{Im}}}
\def\Coker{\mathop{\mathrm{Coker}}}
\def\Coim{\mathop{\mathrm{Coim}}}

% Boxtimes
%
\DeclareMathSymbol{\boxtimes}{\mathbin}{AMSa}{"02}

%
% Macros for moduli stacks/spaces
%
\def\QCohstack{\mathcal{QC}\!\mathit{oh}}
\def\Cohstack{\mathcal{C}\!\mathit{oh}}
\def\Spacesstack{\mathcal{S}\!\mathit{paces}}
\def\Quotfunctor{\mathrm{Quot}}
\def\Hilbfunctor{\mathrm{Hilb}}
\def\Curvesstack{\mathcal{C}\!\mathit{urves}}
\def\Polarizedstack{\mathcal{P}\!\mathit{olarized}}
\def\Complexesstack{\mathcal{C}\!\mathit{omplexes}}
% \Pic is the operator that assigns to X its picard group, usage \Pic(X)
% \Picardstack_{X/B} denotes the Picard stack of X over B
% \Picardfunctor_{X/B} denotes the Picard functor of X over B
\def\Pic{\mathop{\mathrm{Pic}}\nolimits}
\def\Picardstack{\mathcal{P}\!\mathit{ic}}
\def\Picardfunctor{\mathrm{Pic}}
\def\Deformationcategory{\mathcal{D}\!\mathit{ef}}
\usepackage{cleveref}
\crefname{originallemma}{original lemma}{original lemmas}
\Crefname{originallemma}{Original lemma}{Original lemmas}
\makeatletter
\newlabel{lemma-locally-constant-on-connected-geom-unibranch}{{0GIY}{source}{}{sourcepointer-1}{}}
\newlabel{lemma-locally-constant-on-connected-geom-unibranch@cref}{{[originallemma][0][]0GIY}{[1][0][]source}{}{}{}}
\makeatother
\crefname{originallemma}{original lemma}{original lemmas}
\Crefname{originallemma}{Original lemma}{Original lemmas}
\makeatletter
\newlabel{pione-section-finite-etale-under-galois}{{03SF}{source}{}{sourcepointer-2}{}}
\newlabel{pione-section-finite-etale-under-galois@cref}{{[originallemma][0][]03SF}{[1][0][]source}{}{}{}}
\makeatother
\crefname{originalremark}{original remark}{original remarks}
\Crefname{originalremark}{Original remark}{Original remarks}
\makeatletter
\newlabel{remark-functorial-locally-constant-on-connected}{{0DV6}{source}{}{sourcepointer-3}{}}
\newlabel{remark-functorial-locally-constant-on-connected@cref}{{[originalremark][0][]0DV6}{[1][0][]source}{}{}{}}
\makeatother
\crefname{originallemma}{original lemma}{original lemmas}
\Crefname{originallemma}{Original lemma}{Original lemmas}
\makeatletter
\newlabel{lemma-action-l-group}{{0GIZ}{source}{}{sourcepointer-4}{}}
\newlabel{lemma-action-l-group@cref}{{[originallemma][0][]0GIZ}{[1][0][]source}{}{}{}}
\makeatother
\title{Filtering a bounded torsion locally constant module sheaf after a cover of degree prime to l}
\author{The Stacks Project Authors}
\date{stacks-project:etale-cohomology.tex:local-snapshot-20261011; GNU FDL 1.2 or later}
\begin{document}
\maketitle
\noindent\textit{Editorial note.} All delivered mathematical text is the author's own, copied from the pinned Stacks Project chapter file etale-cohomology.tex; no statement or proof step was rewritten, reordered or removed.\par\medskip
\noindent\textit{Editorial note.} The proof is the author's complete argument; the cover, its degree and the filtration are all produced inside the counted span.\par\medskip
\noindent\textit{Editorial note.} Every labelled reference occurring inside the counted spans is carried in assembly.external\_references as a non-load-bearing navigation pointer (same-chapter ones at their definition site, cross-chapter ones at the invocation inside the counted span), so no load-bearing mathematics is replaced by a stub. No bibliography citation occurs inside the counted spans.\par\medskip
\noindent\textit{Editorial note.} The checked chapter list for same-name collisions is the set of stacks chapters staged in this repository: chow, cohomology, divisors, etale-cohomology, flat, more-algebra, more-morphisms, morphisms, spaces-more-morphisms, varieties together with this card's chapter etale-cohomology.\par\medskip
\noindent\textit{Editorial note.} The preamble is the chapter's own native preamble excerpt already used for this source: the Xy-pic setup (native preamble lines 14-18) and the native theorem declarations with the mathematical macros (native preamble lines 154-217).\par\medskip
\noindent\textit{Allowed standard prerequisites (editorial).}\begin{itemize}
\item Classification of finite type locally constant sheaves of Lambda-modules by representations of the fundamental group. Same chapter, the equivalence between finite type locally constant sheaves of Lambda-modules on an irreducible geometrically unibranch scheme and finite Lambda-modules with continuous action of the fundamental group. Supplies the translation of the problem into a statement about modules over a finite group.
\item Structure of modules over the group algebra of a finite group of order prime to l. Same chapter together with the group-theoretic input used by the author: the group algebra of a finite group of order prime to l is a product of matrix algebras over the residue ring, so its modules admit filtrations with successive quotients induced from simple modules of the group. Supplies the filtration produced after passing to the cover.
\end{itemize}\medskip
\ifdefined\hypertarget\hypertarget{sourcepointer-1}{}\fi
\noindent\textit{Source reference.} Original lemma 0GIY: sources/stacks-etale-cohomology/pinned-original-source.tex, source lines 10200--10200. Navigation pointer invoked at proof line 10494; the counted span contains the pointer only and the cited result is not reproduced inside it, so the pointer is registered as a non-load-bearing navigation stub in the established form for Stacks material. No bibliography citation is involved.\par\medskip
\ifdefined\hypertarget\hypertarget{sourcepointer-2}{}\fi
\noindent\textit{Source reference.} Original lemma 03SF: sources/stacks-etale-cohomology/pinned-original-source.tex, source lines 10437--10437. Navigation pointer invoked at proof line 10437; the counted span contains the pointer only and the cited result is not reproduced inside it, so the pointer is registered as a non-load-bearing navigation stub in the established form for Stacks material. No bibliography citation is involved.\par\medskip
\ifdefined\hypertarget\hypertarget{sourcepointer-3}{}\fi
\noindent\textit{Source reference.} Original remark 0DV6: sources/stacks-etale-cohomology/pinned-original-source.tex, source lines 10294--10294. Navigation pointer invoked at proof line 10446; the counted span contains the pointer only and the cited result is not reproduced inside it, so the pointer is registered as a non-load-bearing navigation stub in the established form for Stacks material. No bibliography citation is involved.\par\medskip
\ifdefined\hypertarget\hypertarget{sourcepointer-4}{}\fi
\noindent\textit{Source reference.} Original lemma 0GIZ: sources/stacks-etale-cohomology/pinned-original-source.tex, source lines 10460--10460. Navigation pointer invoked at proof line 10516; the counted span contains the pointer only and the cited result is not reproduced inside it, so the pointer is registered as a non-load-bearing navigation stub in the established form for Stacks material. No bibliography citation is involved.\par\medskip
\begin{lemma}
\label{lemma-pullback-filtered-modules}
Let $S$ be an irreducible, geometrically unibranch scheme.
Let $\ell$ be a prime number and $n \geq 1$. Let $\Lambda$
be a Noetherian ring. Let $\mathcal{F}$ be a finite type,
locally constant sheaf of $\Lambda$-modules on $S_\etale$
which is annihilated by $\ell^n$. Then there exists a
finite \'etale morphism $f : T \to S$ of degree prime to $\ell$
such that $f^{-1}\mathcal{F}$ has a finite filtration whose
successive quotients are of the form $\underline{M}_T$
for some finite $\Lambda$-modules $M$.
\end{lemma}

\begin{proof}
Choose a geometric point $\overline{s}$ of $S$.
Via the equivalence of
Lemma \ref{lemma-locally-constant-on-connected-geom-unibranch}
the sheaf $\mathcal{F}$ corresponds to a finite $\Lambda$-module
$M$ with a continuous $\pi_1(S, \overline{s})$-action.
Let $G \subset \text{Aut}(V)$ be the image of the homomorphism
$\rho : \pi_1(S, \overline{s}) \to \text{Aut}(M)$ giving the action.
Observe that $G$ is finite as $M$ is a finite $\Lambda$-module
(see proof of Lemma \ref{lemma-locally-constant-on-connected-geom-unibranch}).
The surjective continuous homomorphism
$\overline{\rho} : \pi_1(S, \overline{s}) \to G$
corresponds to a Galois object $Y \to S$ of
$\textit{F\'Et}_S$ with automorphism group $G = \text{Aut}(Y/S)$, see
Fundamental Groups, Section \ref{pione-section-finite-etale-under-galois}.
Let $H \subset G$ be an $\ell$-Sylow subgroup.
We claim that $T = Y/H \to S$ works. Namely, let $\overline{t} \in T$
be a geometric point over $\overline{s}$. The image of
$\pi_1(T, \overline{t}) \to \pi_1(S, \overline{s})$
is $(\overline{\rho})^{-1}(H)$ as follows from the functorial
nature of fundamental groups. Hence the action of $\pi_1(T, \overline{t})$
on $M$ corresponding to $f^{-1}\mathcal{F}$ is through
the map $\pi_1(T, \overline{t}) \to H$, see
Remark \ref{remark-functorial-locally-constant-on-connected}.
Let $0 = M_0 \subset M_1 \subset \ldots \subset M_t = M$
be as in Lemma \ref{lemma-action-l-group}.
This induces a filtration
$0 = \mathcal{F}_0 \subset \mathcal{F}_1 \subset \ldots
\subset \mathcal{F}_t = f^{-1}\mathcal{F}$
such that the successive quotients are constant with
value $M_{i + 1}/M_i$.
Finally, the degree of $T = Y/H \to S$ is prime to $\ell$
as it is equal to the index of $H$ in $G$.
\end{proof}

\end{document}
